\chapter{Abstract Data Types in C}

\textbf{Readings:} CP:AMA 19.3--19.5.

In Unit~03 we introduced ADTs as collections described only by their operations. Now that we
have opaque structures, dynamic memory and linked data, we can implement them properly in C and
compare implementations by efficiency.

\section{Designing an ADT}

Designing an ADT means deciding:
\begin{enumerate}
  \item \textbf{The operations} and their exact behaviour (documented in the \code{.h} file).
  \item \textbf{Error handling:} what happens on invalid use? Either make it a precondition
    (``requires: stack is not empty'', checked with \code{assert}) or report it (return a
    \code{bool}, return \code{NULL}, a special value).
  \item \textbf{Memory ownership:} who allocates and frees what? Does the ADT store a
    \emph{copy} of the client's data (e.g.\ duplicate a string) or just the pointer it was
    given? Who frees items when the ADT is destroyed?
  \item \textbf{The implementation}: the data structure and the \emph{invariants} it maintains
    (e.g.\ ``the array is sorted'', ``\code{len <= maxlen}'', ``BST ordering holds''), chosen for
    the efficiency the clients need.
\end{enumerate}

\begin{important}[Invariants]
A \emph{representation invariant} is a property of the internal data that every interface
function may assume on entry and must restore before returning. Because the structure is
opaque, only the module's own functions can break an invariant --- this is exactly why
information hiding gives \emph{security}.
\end{important}

\section{The Dictionary ADT}

\subsection{Interface}

We implement a dictionary from \code{int} keys to string values. The dictionary stores its own
copies of the strings.
\begin{lstlisting}[style=c]
// dictionary.h: a dictionary ADT (int keys -> string values)
#include <stdbool.h>

struct dictionary;

// dict_create() creates a new, empty dictionary
// effects: allocates memory (client must call dict_destroy)
struct dictionary *dict_create(void);

// dict_destroy(d) frees all memory used by d (including stored strings)
void dict_destroy(struct dictionary *d);

// dict_insert(key, val, d) adds key -> val to d; if key is already
//   present, its value is replaced. d stores a copy of val.
// effects: modifies d
// time: O(h + m), h is the height of the tree, m is the length of val
void dict_insert(int key, const char *val, struct dictionary *d);

// dict_lookup(key, d) produces the value for key, or NULL if not present
//   the produced string must not be modified or freed by the client
// time: O(h)
const char *dict_lookup(int key, const struct dictionary *d);

// dict_remove(key, d) removes key from d (no effect if not present)
// effects: modifies d
// time: O(h)
void dict_remove(int key, struct dictionary *d);
\end{lstlisting}

\subsection{Implementation with a BST}

\begin{lstlisting}[style=c]
// dictionary.c
#include <assert.h>
#include <stdlib.h>
#include <string.h>
#include "dictionary.h"

struct bstnode {
  int key;
  char *val;                     // owned by the dictionary
  struct bstnode *left;
  struct bstnode *right;
};

struct dictionary {
  struct bstnode *root;
};

// my_strdup(s) produces a new heap copy of s
static char *my_strdup(const char *s) {
  char *copy = malloc((strlen(s) + 1) * sizeof(char));   // +1 for '\0'
  strcpy(copy, s);
  return copy;
}

struct dictionary *dict_create(void) {
  struct dictionary *d = malloc(sizeof(struct dictionary));
  d->root = NULL;
  return d;
}

static void free_node(struct bstnode *node) {
  if (node) {
    free_node(node->left);
    free_node(node->right);
    free(node->val);
    free(node);
  }
}

void dict_destroy(struct dictionary *d) {
  assert(d);
  free_node(d->root);
  free(d);
}

static struct bstnode *insert_node(int key, const char *val,
                                   struct bstnode *node) {
  if (node == NULL) {
    struct bstnode *leaf = malloc(sizeof(struct bstnode));
    leaf->key = key;
    leaf->val = my_strdup(val);
    leaf->left = NULL;
    leaf->right = NULL;
    return leaf;
  }
  if (key < node->key) {
    node->left = insert_node(key, val, node->left);
  } else if (key > node->key) {
    node->right = insert_node(key, val, node->right);
  } else {                                    // replace existing value
    free(node->val);
    node->val = my_strdup(val);
  }
  return node;
}

void dict_insert(int key, const char *val, struct dictionary *d) {
  assert(d && val);
  d->root = insert_node(key, val, d->root);
}

const char *dict_lookup(int key, const struct dictionary *d) {
  assert(d);
  const struct bstnode *node = d->root;
  while (node) {
    if (key == node->key) return node->val;
    node = (key < node->key) ? node->left : node->right;
  }
  return NULL;
}
\end{lstlisting}

\subsection{Removing from a BST}

Removing a node has three cases:
\begin{enumerate}
  \item \textbf{No children:} just free it (the parent's pointer becomes \code{NULL}).
  \item \textbf{One child:} replace the node with its child.
  \item \textbf{Two children:} find the \emph{successor} (the smallest key in the right
    subtree, i.e.\ the leftmost node there), copy its key and value into the node, and then
    remove the successor from the right subtree (it has no left child, so that is case 1 or 2).
\end{enumerate}
\begin{lstlisting}[style=c]
// remove_node(key, node) removes key from the subtree rooted at node
//   and produces the new root of that subtree
static struct bstnode *remove_node(int key, struct bstnode *node) {
  if (node == NULL) return NULL;            // key not found
  if (key < node->key) {
    node->left = remove_node(key, node->left);
  } else if (key > node->key) {
    node->right = remove_node(key, node->right);
  } else if (node->left == NULL || node->right == NULL) {
    struct bstnode *child = node->left ? node->left : node->right;
    free(node->val);
    free(node);
    return child;                           // cases 1 and 2
  } else {
    struct bstnode *succ = node->right;     // case 3
    while (succ->left) {
      succ = succ->left;
    }
    free(node->val);
    node->key = succ->key;
    node->val = my_strdup(succ->val);
    node->right = remove_node(succ->key, node->right);
  }
  return node;
}

void dict_remove(int key, struct dictionary *d) {
  assert(d);
  d->root = remove_node(key, d->root);
}
\end{lstlisting}
\begin{aside}
Instead of copying the successor's string and then freeing the original, we could move the
pointer over and arrange for the recursive call not to free it. The version above is simpler to
reason about.
\end{aside}

\subsection{Comparing dictionary implementations}

\begin{center}
\begin{tabular}{lccc}
  \toprule
  Implementation & lookup & insert & remove \\
  \midrule
  unsorted array / linked list & $\bigO(n)$ & $\bigO(n)$\textsuperscript{*} & $\bigO(n)$ \\
  sorted array & $\bigO(\log n)$ & $\bigO(n)$ & $\bigO(n)$ \\
  BST (worst case) & $\bigO(n)$ & $\bigO(n)$ & $\bigO(n)$ \\
  balanced BST (AVL, CS 240) & $\bigO(\log n)$ & $\bigO(\log n)$ & $\bigO(\log n)$ \\
  hash table (CS 240) & $\bigO(1)$ average & $\bigO(1)$ average & $\bigO(1)$ average \\
  \bottomrule
\end{tabular}
\end{center}
\textsuperscript{*}$\bigO(1)$ if duplicates are not checked (but then insert must search to
replace a value).

The point of the ADT is that a client of \code{dictionary.h} does not change at all if we swap
one implementation for another.

\section{Stack, Queue and Sequence Implementations}

\subsection{Stack}

Both a dynamic array (Unit~10) and a linked list (push/pop at the front) give $\bigO(1)$
operations (amortized for the array). The array version has better memory locality; the list
version never needs to copy.

\subsection{Queue}

A linked list with front and back pointers (Unit~11) gives $\bigO(1)$ add-back and
remove-front. An array-based queue would need $\bigO(n)$ to remove from the front by shifting,
unless we use a \emph{circular buffer}:
\begin{lstlisting}[style=c]
struct queue {
  int *data;
  int maxlen;
  int front;     // index of the front item
  int len;       // number of items
};

// time: O(1) amortized
void queue_add_back(int item, struct queue *q) {
  if (q->len == q->maxlen) {
    int *new_data = malloc(2 * q->maxlen * sizeof(int));
    for (int i = 0; i < q->len; ++i) {        // unwrap into the new array
      new_data[i] = q->data[(q->front + i) % q->maxlen];
    }
    free(q->data);
    q->data = new_data;
    q->front = 0;
    q->maxlen *= 2;
  }
  q->data[(q->front + q->len) % q->maxlen] = item;
  ++q->len;
}

// requires: q is not empty
// time: O(1)
int queue_remove_front(struct queue *q) {
  assert(q->len > 0);
  int item = q->data[q->front];
  q->front = (q->front + 1) % q->maxlen;
  --q->len;
  return item;
}
\end{lstlisting}
The indices ``wrap around'' using \code{\%}, so the items occupy a contiguous range modulo the
capacity.

\subsection{Sequence}

\begin{center}
\begin{tabular}{lccc}
  \toprule
  Implementation & item\_at & insert\_at / remove\_at & length \\
  \midrule
  dynamic array & $\bigO(1)$ & $\bigO(n)$ & $\bigO(1)$ \\
  linked list (with length) & $\bigO(i)$ & $\bigO(i)$ & $\bigO(1)$ \\
  balanced tree with subtree sizes & $\bigO(\log n)$ & $\bigO(\log n)$ & $\bigO(1)$ \\
  \bottomrule
\end{tabular}
\end{center}
No implementation is best for every operation; the right choice depends on which operations the
client uses most.

\section{Set ADT}

A \emph{set} stores unique items with operations \code{member}, \code{add}, \code{remove}, and
often \code{union}, \code{intersection}, \code{size}. It is essentially a dictionary with keys but
no values, so the same implementations apply. With sorted arrays or sorted lists, union and
intersection can be done in $\bigO(n + m)$ by a merge-like simultaneous walk.

\section{Iterating Over an ADT}

Clients often need to visit every item, but they cannot access the internal structure. Common
solutions:
\begin{itemize}
  \item an interface function that produces a new array of all items (in order), which the client
    must free;
  \item a ``map''-style function that applies a client's function pointer to every item;
  \item an iterator interface (\code{first}, \code{next}, \code{done}), as in C++ (CS 246).
\end{itemize}
\begin{lstlisting}[style=c]
// dict_keys(d, len) produces a new array of the keys of d in increasing
//   order and stores the number of keys in *len
// effects: allocates memory (caller must free), modifies *len
// time: O(n)
int *dict_keys(const struct dictionary *d, int *len);
\end{lstlisting}

\section{Summary}
\begin{itemize}
  \item ADT design: operations, error handling (requires vs.\ return codes), memory ownership,
    and implementation invariants.
  \item Dictionary via BST: lookup/insert/remove are $\bigO(h)$; removal with two children uses
    the successor.
  \item Compare implementations in a table; choose based on which operations matter.
  \item Queues can use a linked list or a circular array; sequences trade indexing against
    insertion.
  \item Clients can iterate through an exported array, a map with function pointers, or an
    iterator.
\end{itemize}
